% Engendré par scripts/collect_figures.py — NE PAS ÉDITER À LA MAIN.
% 40 figures.

{\setlength{\tabcolsep}{4pt}
\begin{longtable}{@{}>{\raggedright\arraybackslash}p{0.18\textwidth}>{\raggedright\arraybackslash}p{0.27\textwidth}>{\raggedright\arraybackslash}p{0.45\textwidth}@{}}
\toprule
\textbf{figure} & \textbf{provenance} & \textbf{ce qu'elle établit} \\
\midrule
\endhead
{\small\texttt{m01\_\allowbreak{}reseau\_\allowbreak{}credit}}\newline{\scriptsize\texttt{daa5f01fe9a4}} & Simulation Lab, run 20260718\_\allowbreak{}033746\_\allowbreak{}8b0c5acb (M4B, point central, graine 14) & Le réseau de crédit final : aucun type déclaré, et 182 entités sur 228 sont simultanément prêteuses et emprunteuses. \\
{\small\texttt{m02\_\allowbreak{}vue\_\allowbreak{}macro}}\newline{\scriptsize\texttt{2e48959eb0d1}} & Simulation Lab, run 20260718\_\allowbreak{}033746\_\allowbreak{}8b0c5acb & Capital, crédit nouveau, population et levier sur 4000 pas : le système est borné et stationnaire, sans dérive ni extinction. \\
{\small\texttt{t01\_\allowbreak{}artefact\_\allowbreak{}cohorte}}\newline{\scriptsize\texttt{e093fd184fac}} & recherche/\allowbreak{}analyse\_\allowbreak{}distributions\_\allowbreak{}taille\_\allowbreak{}revenu, campagne du WIP du 27 avril (163 runs stationnaires sur 805) & Le résultat négatif fondateur : la bimodalité lue comme « banques et travailleurs » est une structure d'âge, pas deux classes économiques. \\
{\small\texttt{t02\_\allowbreak{}derive\_\allowbreak{}queue}}\newline{\scriptsize\texttt{4408f6ce8620}} & même campagne & La première queue lourde n'était pas stationnaire : son exposant dérive avec l'horizon d'observation. \\
{\small\texttt{t03\_\allowbreak{}branchement\_\allowbreak{}m4}}\newline{\scriptsize\texttt{7a57ca8f11c0}} & m4\_\allowbreak{}credit\_\allowbreak{}soc\_\allowbreak{}fable, rapport 01\_\allowbreak{}soc\_\allowbreak{}final & M4 : le rapport de branchement se stabilise vers 0,30 une fois la règle annulation–destruction en place. \\
{\small\texttt{t04\_\allowbreak{}taille\_\allowbreak{}finie\_\allowbreak{}m4}}\newline{\scriptsize\texttt{26a7595238c4}} & m4\_\allowbreak{}credit\_\allowbreak{}soc\_\allowbreak{}fable, rapport 01\_\allowbreak{}soc\_\allowbreak{}final & M4 : la coupure des avalanches croît avec la taille du système — signature de taille finie, condition nécessaire d'un candidat SOC. \\
{\small\texttt{t05\_\allowbreak{}effet\_\allowbreak{}echelle\_\allowbreak{}m42}}\newline{\scriptsize\texttt{dcdd296f22ed}} & m4\_\allowbreak{}2\_\allowbreak{}credit\_\allowbreak{}soc, campagne M4.2 & M4.2 : l'effet attribué naïvement à la concavité γ est dominé par un déplacement d'échelle du capital. \\
{\small\texttt{p01\_\allowbreak{}lorenz}}\newline{\scriptsize\texttt{39a211d47ecf}} & Simulation Lab, run 20260718\_\allowbreak{}033746\_\allowbreak{}8b0c5acb & L'inégalité ne vit pas dans la taille productive (G = 0,215) mais dans le revenu d'intérêt (G = 0,595). \\
{\small\texttt{p02\_\allowbreak{}structure\_\allowbreak{}nw}}\newline{\scriptsize\texttt{a26543aae3ee}} & recherche/\allowbreak{}sensibilite\_\allowbreak{}m4b, campagne de 594 simulations & Structure de la valeur nette : cycle de vie débitrice puis créancière, et trou d'insolvabilité au décès. \\
{\small\texttt{p03\_\allowbreak{}vie\_\allowbreak{}par\_\allowbreak{}taille}}\newline{\scriptsize\texttt{8887ab830575}} & Simulation Lab, run 20260718\_\allowbreak{}033746\_\allowbreak{}8b0c5acb & L'espérance de vie restante croît avec la taille : les grandes entités survivent plus longtemps, sans être immortelles. \\
{\small\texttt{p04\_\allowbreak{}avalanches\_\allowbreak{}structure}}\newline{\scriptsize\texttt{4f7ede37f1bf}} & Simulation Lab, run 20260718\_\allowbreak{}033746\_\allowbreak{}8b0c5acb & Structure causale des avalanches : la part des racines décroît avec la taille et la profondeur atteint 10 — c'est de la propagation, pas de la synchronisation. \\
{\small\texttt{p05\_\allowbreak{}cascades\_\allowbreak{}rank\_\allowbreak{}size}}\newline{\scriptsize\texttt{95e15b6fec17}} & Simulation Lab, run 20260718\_\allowbreak{}033746\_\allowbreak{}8b0c5acb & Densité, CCDF et volume des cascades : loi de puissance tronquée, coupure visible dans la fenêtre observable. \\
{\small\texttt{p06\_\allowbreak{}scaling\_\allowbreak{}avalanches}}\newline{\scriptsize\texttt{b983649476fb}} & recherche/\allowbreak{}sensibilite\_\allowbreak{}m4b, phase de confirmation & Les lois d'échelle des avalanches en fonction de la taille du système : $s_{\max} \propto \lambda^{0,58}$ et $\chi \propto \lambda^{0,51}$. \\
{\small\texttt{p07\_\allowbreak{}hierarchie\_\allowbreak{}prcc}}\newline{\scriptsize\texttt{e27175372621}} & recherche/\allowbreak{}sensibilite\_\allowbreak{}m4b, criblage LHS & Hiérarchie des paramètres : $\sigma$, $k$, $K_0$, $\delta$ et $\lambda$ ont des rôles distincts et séparables. \\
{\small\texttt{p08\_\allowbreak{}gini\_\allowbreak{}renouvellement}}\newline{\scriptsize\texttt{99507cfc5529}} & recherche/\allowbreak{}sensibilite\_\allowbreak{}m4b, phase de confirmation & Gini et renouvellement sur graines de confirmation : les deux inégalités, capital et valeur nette, répondent parfois en sens opposé. \\
{\small\texttt{p09\_\allowbreak{}cycles\_\allowbreak{}structure}}\newline{\scriptsize\texttt{5f8d465986c5}} & recherche/\allowbreak{}sensibilite\_\allowbreak{}m4b, extension cycles (post-audit) & Structure des épisodes d'activité : courts, asymétriques, sans longue récession endogène. \\
{\small\texttt{p10\_\allowbreak{}cycles\_\allowbreak{}sensibilite}}\newline{\scriptsize\texttt{03e540350fef}} & recherche/\allowbreak{}sensibilite\_\allowbreak{}m4b, extension cycles (post-audit) & Ce qui pilote les épisodes : l'amplitude suit σ et K₀, la fréquence n'est expliquée de façon robuste par aucun paramètre. \\
{\small\texttt{p11\_\allowbreak{}renouvellement\_\allowbreak{}top}}\newline{\scriptsize\texttt{1fb0a1f62b71}} & Simulation Lab, run 20260718\_\allowbreak{}033746\_\allowbreak{}8b0c5acb & Le décile supérieur se renouvelle intégralement en moins de cent pas : c'est un renouvellement démographique, et non une mobilité sociale comparable aux données. \\
{\small\texttt{p12\_\allowbreak{}revenu\_\allowbreak{}empirique}}\newline{\scriptsize\texttt{daec7099081e}} & Simulation Lab, run 20260718\_\allowbreak{}033746\_\allowbreak{}8b0c5acb & Le revenu brut du modèle, sans ajustement imposé : ni corps exponentiel net, ni queue de Pareto convaincante sur cette observable. \\
{\small\texttt{p13\_\allowbreak{}queue\_\allowbreak{}interets}}\newline{\scriptsize\texttt{3d7090de2469}} & m4\_\allowbreak{}2b\_\allowbreak{}credit\_\allowbreak{}soc, campagne M4.2B & La queue d'intérêts telle qu'elle se voit sur les graphiques — l'observation qui a motivé le test de Pareto. \\
{\small\texttt{p14\_\allowbreak{}erreur\_\allowbreak{}seuil}}\newline{\scriptsize\texttt{7eae899c2345}} & m4\_\allowbreak{}2b\_\allowbreak{}credit\_\allowbreak{}soc, campagne M4.2B & Le contrôle d'autosimilarité mal orienté, puis corrigé : à seuil×2 la queue médiane tombe de 113 à 17 points. L'existence de Pareto reste une hypothèse. \\
{\small\texttt{r01\_\allowbreak{}reponse\_\allowbreak{}prod}}\newline{\scriptsize\texttt{13fa2cce2157}} & m4\_\allowbreak{}3live\_\allowbreak{}credit\_\allowbreak{}soc, campagne M4.3Live & La mesure centrale du rebond : à portée partielle la réponse est forte puis s'éteint avec la cohorte traitée ; à portée globale elle s'établit à +36 \% pour A×1,5, donc sous-proportionnelle. \\
{\small\texttt{r02\_\allowbreak{}ablation\_\allowbreak{}k0}}\newline{\scriptsize\texttt{3841b13b9646}} & m4\_\allowbreak{}3live\_\allowbreak{}credit\_\allowbreak{}soc, campagne M4.3Live & Le pivot du verdict : selon que la dotation de naissance K₀ suit ou non l'échelle technologique, la population se contracte ou non. \\
{\small\texttt{r03\_\allowbreak{}covariance\_\allowbreak{}echelle}}\newline{\scriptsize\texttt{90a37fc95450}} & m4\_\allowbreak{}3live\_\allowbreak{}credit\_\allowbreak{}soc, campagne M4.3Live & La covariance d'échelle du moteur, vérifiée numériquement : sous compensation de K₀, l'élasticité vaut exactement 1/(1−γ). \\
{\small\texttt{f01\_\allowbreak{}ancrages}}\newline{\scriptsize\texttt{380d0c4ce946}} & m4\_\allowbreak{}4\_\allowbreak{}rebond\_\allowbreak{}credit\_\allowbreak{}soc, engendrée par scripts/\allowbreak{}make\_\allowbreak{}figures.py & Les trois élasticités de la lignée sont reproduites sur une campagne neuve. \\
{\small\texttt{f05\_\allowbreak{}position\_\allowbreak{}nette}}\newline{\scriptsize\texttt{2a343fbc63d3}} & m4\_\allowbreak{}4\_\allowbreak{}rebond\_\allowbreak{}credit\_\allowbreak{}soc, engendrée par scripts/\allowbreak{}make\_\allowbreak{}figures.py & La position nette de crédit par âge, et son déplacement sous intervention. \\
{\small\texttt{f06\_\allowbreak{}mortalite\_\allowbreak{}rotation}}\newline{\scriptsize\texttt{d68a937c47b2}} & m4\_\allowbreak{}4\_\allowbreak{}rebond\_\allowbreak{}credit\_\allowbreak{}soc, engendrée par scripts/\allowbreak{}make\_\allowbreak{}figures.py & L'exposant qui relie mortalité et rotation du crédit n'est pas universel : il dépend du levier. \\
{\small\texttt{f07\_\allowbreak{}chaine}}\newline{\scriptsize\texttt{fc2e1cb4c39e}} & m4\_\allowbreak{}4\_\allowbreak{}rebond\_\allowbreak{}credit\_\allowbreak{}soc, engendrée par scripts/\allowbreak{}make\_\allowbreak{}figures.py & La chaîne causale de la contraction, maillon par maillon, avec la valeur de chaque élasticité. \\
{\small\texttt{f08\_\allowbreak{}avalanches\_\allowbreak{}ccdf}}\newline{\scriptsize\texttt{3e58411cb80c}} & m4\_\allowbreak{}4\_\allowbreak{}rebond\_\allowbreak{}credit\_\allowbreak{}soc, engendrée par scripts/\allowbreak{}make\_\allowbreak{}figures.py & Les CCDF d'avalanches par bras : la forme survit au changement de productivité. \\
{\small\texttt{f09\_\allowbreak{}b1\_\allowbreak{}b2}}\newline{\scriptsize\texttt{3a3b08fdc4ac}} & m4\_\allowbreak{}4\_\allowbreak{}rebond\_\allowbreak{}credit\_\allowbreak{}soc, engendrée par scripts/\allowbreak{}make\_\allowbreak{}figures.py & Les deux estimateurs de branchement, et pourquoi il en faut deux. \\
{\small\texttt{f10\_\allowbreak{}k0}}\newline{\scriptsize\texttt{2c1316cb7b87}} & m4\_\allowbreak{}4\_\allowbreak{}rebond\_\allowbreak{}credit\_\allowbreak{}soc, engendrée par scripts/\allowbreak{}make\_\allowbreak{}figures.py & La fragilité attribuée au rebond n'est pas la sienne : 97,9 \% de l'effet est annulé par la compensation de dotation. \\
{\small\texttt{f12\_\allowbreak{}alpha\_\allowbreak{}bras}}\newline{\scriptsize\texttt{62c50e375bc3}} & m4\_\allowbreak{}4\_\allowbreak{}rebond\_\allowbreak{}credit\_\allowbreak{}soc, engendrée par scripts/\allowbreak{}make\_\allowbreak{}figures.py & Les exposants de queue ne bougent pas d'un bras à l'autre. \\
{\small\texttt{f11\_\allowbreak{}couverture}}\newline{\scriptsize\texttt{c1408a137036}} & m4\_\allowbreak{}4\_\allowbreak{}rebond\_\allowbreak{}credit\_\allowbreak{}soc, engendrée par scripts/\allowbreak{}make\_\allowbreak{}figures.py & La couverture de queue n'est plus l'obstacle, et l'exposant ne bouge pas quand l'échantillon grandit de moitié. \\
{\small\texttt{f16\_\allowbreak{}bargain\_\allowbreak{}systeme}}\newline{\scriptsize\texttt{701670760a09}} & m4\_\allowbreak{}4\_\allowbreak{}rebond\_\allowbreak{}credit\_\allowbreak{}soc, engendrée par scripts/\allowbreak{}make\_\allowbreak{}figures.py & Ce que le partage déplace : branchement, tension, rotation et taux du carnet. \\
{\small\texttt{f18\_\allowbreak{}rotation\_\allowbreak{}rho}}\newline{\scriptsize\texttt{f8e3cc4e7bfe}} & m4\_\allowbreak{}4\_\allowbreak{}rebond\_\allowbreak{}credit\_\allowbreak{}soc, engendrée par scripts/\allowbreak{}make\_\allowbreak{}figures.py & Une relation héritée qui ne survit pas au changement de levier. \\
{\small\texttt{f19\_\allowbreak{}ablation}}\newline{\scriptsize\texttt{79aee50ff832}} & m4\_\allowbreak{}4\_\allowbreak{}rebond\_\allowbreak{}credit\_\allowbreak{}soc, engendrée par scripts/\allowbreak{}make\_\allowbreak{}figures.py & Ce qui explique l'écart de branchement avec la lignée M4B : 87,8 \% attribués, un huitième résiduel. \\
{\small\texttt{f20\_\allowbreak{}rampes}}\newline{\scriptsize\texttt{ecb7337f793c}} & m4\_\allowbreak{}4\_\allowbreak{}rebond\_\allowbreak{}credit\_\allowbreak{}soc, engendrée par scripts/\allowbreak{}make\_\allowbreak{}figures.py & Montée et descente du partage se superposent : retard de relaxation, pas dépendance au chemin. \\
{\small\texttt{f21\_\allowbreak{}hasard}}\newline{\scriptsize\texttt{e4d5f553d85a}} & m4\_\allowbreak{}4\_\allowbreak{}rebond\_\allowbreak{}credit\_\allowbreak{}soc, engendrée par scripts/\allowbreak{}make\_\allowbreak{}figures.py & La convexité posée ne tient pas : le risque n'est pas convexe dans le fardeau, et la courbure meurt au conditionnement. \\
{\small\texttt{f22\_\allowbreak{}ponderation}}\newline{\scriptsize\texttt{c5db42797aab}} & m4\_\allowbreak{}4\_\allowbreak{}rebond\_\allowbreak{}credit\_\allowbreak{}soc, engendrée par scripts/\allowbreak{}make\_\allowbreak{}figures.py & Le mécanisme, et c'est une identité : 0,53 par contrat, 0,94 par joule, l'écart est Cov(p,Δ)/E[Δ]. \\
{\small\texttt{g03\_\allowbreak{}partage\_\allowbreak{}analytique}}\newline{\scriptsize\texttt{692f06393086}} & m4\_\allowbreak{}4\_\allowbreak{}rebond\_\allowbreak{}credit\_\allowbreak{}soc, engendrée par scripts/\allowbreak{}make\_\allowbreak{}figures.py & Le partage dérivé des seules fonctions du moteur : ½ à la limite infinitésimale, au-delà de 1 sur les paires très inégales. \\
\bottomrule
\end{longtable}}
